A partir de um ensaio ao degrau, com entrada de amplitude $50^\circ\mathrm{C}$
aplicada ao aquecedor~1, identificou-se a dinâmica da planta. A realimentação
baseia-se no sensor~1, enquanto o objetivo de seguimento refere-se à temperatura
relativa registrada pelo sensor~2. Adotou-se a abordagem de Ziegler-Nichols por
ensaio ao degrau com um modelo \gls{fopdt}, obtendo-se uma aproximação para cada
sensor, conforme a \Eqref{planta-sensor-1-zn,planta-sensor-2-zn}:

\begin{align}
\eqlabel{planta-sensor-1-zn}
G_{1}(s) &= \frac{0.92}{188.03s + 1} \cdot e^{-6.37s}\\
\eqlabel{planta-sensor-2-zn}
G_{2}(s) &= \frac{0.91}{183.94s + 1} \cdot e^{-11.66s}
\end{align}

O modelo do sensor~1, usado no projeto, foi discretizado por \gls{zoh} com
$T_s = \SI{15}{\second}$, resultando na \Eqref{planta-gz-zoh}.